% 完整翻译 source-reference/TERMINAL_INDEX.tex, 源行 1--254.
% BEGIN inputs/TERMINAL_INDEX/index/tex/SGA3_TomeIII_Index_body.tex
\chapter*{索引——第三卷}
\addcontentsline{toc}{chapter}{索引——第三卷}
\markboth{索引——第三卷}{索引——第三卷}
\noindent\emph{页码依照 Polo--Gille 重订版.}
\medskip
\begin{multicols}{2}
\small
\setlength{\parskip}{1.2pt}
\interlinepenalty=10000
% BEGIN inputs/TERMINAL_INDEX/index/tex/components/00_tomeIII_index_entries.tex
\Idx{\(\operatorname{ad}(G)\)}{127}
\Idx{伴随（典范分解）}{243}
\Idx{伴随 \(S\)-约化群}{127}
\Idx{伴随根数据}{91}
\Idx{\(\operatorname{ad}(\mathscr R)\)}{95, 98}
\Idx{\(\alpha^*\)}{110}
\Idx{\(\mathfrak g^\alpha\) 与 \(\mathfrak g^{-\alpha}\) 的配对截面}{110}
\Idx{\(A_S(\mathscr R)\)}{223}
\Idx{\(\underline{\operatorname{Aut}}_{S\text{-}\mathrm{gr}}(G)\)}{217}
\Idx{约化群的自同构}{216}
\Idx{根数据的自同构}{98}
\Idx{约化群的 Borel 子群的自同构}{248}
\Idx{\(\operatorname{Aut}(\mathscr R)\)}{98}
\Idx{\(\operatorname{Aut}^s(\mathscr R)\)}{99}
\Idx{\(B^{\mathrm{ad}}\)}{224}
\Idx{Borel 子群}{131}
\Idx{\(\underline{\operatorname{Bor}}(G)\)}{156}
\Idx{Bruhat 分解}{149, 151}
\Idx{Cartan 子群}{128}
\Idx{\(S\)-约化群的中心}{119}
\Idx{Weyl 室}{83}
\Idx{Chevalley 群概形}{209}
\Idx{\(\operatorname{Cl}_{\mathrm{crit}}\)}{292}
\Idx{根的闭集}{71}
\Idx{约化群概形的余根}{110, 111}
\Idx{根数据的余根}{64}
\Idx{无穷小余根 \(H_\alpha\)}{114}
\Idx{\(\operatorname{corad}(G)\)}{171}
\Idx{余根基环面}{171}
\Idx{根数据的余根基}{92}
\Idx{\(\operatorname{corad}(\mathscr R)\)}{92}
\Idx{Killing 对}{135}
\Idx{\(\underline{\operatorname{Crit}}(G)\)}{292}
\Idx{临界子群}{292}
\Idx{\(C\)-临界环面}{292}
\Idx{\(C^i(T)\)}{262}
\Idx{\(\underline{\mathrm{CT}}\)}{293}
\Idx{约化群的分裂结构}{113}
\Idx{分裂约化群}{113}
\Idx{分裂环面}{10}
\Idx{可分裂约化群}{113}
\Idx{\(\operatorname{der}(\mathscr R)\)}{95}
\Idx{根数据的 Dynkin 图}{105}
\Idx{根数据}{63}
\Idx{伴随根数据}{91}
\Idx{对偶根数据}{64}
\Idx{诱导或余诱导根数据}{94}
\Idx{不可约根数据}{102}
\Idx{约化根数据}{68}
\Idx{相对根数据}{325}
\Idx{半单根数据}{64}
\Idx{单连通根数据}{91}
\Idx{扭根数据}{112}
\Idx{平凡根数据}{64}
\Idx{带基根数据}{179}
\Idx{\(\underline{\operatorname{Dyn}}(G)\)}{228}
\Idx{\(\underline{\operatorname{Dyn}}_0(G)\)}{244}
\Idx{\(\underline{\operatorname{Dyn}}_{0,t}(G)\)}{244}
\Idx{\(E_\Delta(\mathscr R)\)}{98}
\Idx{\(E^s_\Delta(\mathscr R)\)}{99}
\Idx{当 \(S\) 为局部时, \(G(S)\) 由 \(U_\alpha(S)\) 生成}{151}
\Idx{钉扎}{177}
\Idx{钉扎 \(S\)-约化群}{178}
\Idx{本质自由的 \(S\)-概形}{132, 152}
\Idx{\(\exp\)}{29}
\Idx{\(\exp_\alpha\)}{109}
\Idx{指数映射}{29}
\Idx{\((\mathrm F)\), 公式}{35}
\Idx{\(\mathscr F^\alpha\)}{12}
\Idx{\(\underline{\operatorname{Fib}}(S,G)\)}{215}
\Idx{\(G\) 在 \(S\) 上的形式}{222}
\Idx{\(\mathfrak g^\alpha\)}{5}
\Idx{\(\Gamma_0(\mathscr R)\)}{64}
\Idx{钉扎群的生成元与关系}{182}
\Idx{\(\underline{\operatorname{Gen}}(G)\)}{299}
\Idx{\(\underline{\operatorname{Gen}}(P/Q)\)}{299}
\Idx{\(\underline{\operatorname{Gen}}(/Q)\)}{299}
\Idx{\(G^{q\text{-}\text{ép}.}\)}{233}
\Idx{大胞腔}{118, 271}
\Idx{钉扎群}{216}
\Idx{\textup{(RA)} 型群}{130}
\Idx{\textup{(RR)} 型群}{128}
\Idx{\(S\)-约化群的导出群}{171}
\Idx{拟钉扎群}{230}
\Idx{扩充 Weyl 群}{212}
\Idx{环面的 Weyl 群}{6}
\Idx{根数据的 Weyl 群}{65}
\Idx{根数据的 Weyl 群的生成元与关系}{86}
\Idx{\((G,T)/S\) 的 Weyl 群}{116}
\Idx{半单秩为一时的 Weyl 群}{44}
\Idx{代数闭域上的约化群}{4}
\Idx{任意基上的约化群}{12}
\Idx{与 \(S\)-约化群相联系的半单群}{127}
\Idx{代数闭域上的半单群}{5}
\Idx{任意基上的半单群}{12}
\Idx{\(G_S^{\text{Ép}}(\mathscr R)\)}{209}
\Idx{不可除根}{68}
\Idx{\(S\)-约化群的中心同源}{124}
\Idx{\(S\)-约化群的同源}{124}
\Idx{根数据的同源}{90}
\Idx{\(\underline{\operatorname{Isom}}^{\mathrm{ext}}(G,G')\)}{220}
\Idx{\(\underline{\operatorname{Isom}}^{\mathrm{int},u}(G,G')\)}{226}
\Idx{局部等平凡}{237}
\Idx{半局部等平凡}{237}
\Idx{等型分量}{243}
\Idx{\(\underline{\operatorname{Kil}}(G)\)}{156}
\Idx{\(\Lambda(\mathscr R)\)}{97}
\Idx{\(\underline{\operatorname{Lev}}(P)\)}{282}
\Idx{根数据的 Cartan 矩阵}{104}
\Idx{根数据的态射}{90}
\Idx{钉扎群的态射}{178}
\Idx{分裂群的态射}{122}
\Idx{\(\underline{\operatorname{Opp}}(B)\)}{158}
\Idx{\(\underline{\operatorname{Opp}}(G)\)}{305}
\Idx{\(\underline{\operatorname{Opp}}(/P)\)}{305}
\Idx{相反的抛物子群}{303}
\Idx{\(\operatorname{ord}_\Delta(\alpha)\)}{76}
\Idx{\(\operatorname{Of}(E)\)}{289}
\Idx{\(p_\alpha\)}{27, 109}
\Idx{\(\underline{\operatorname{Par}}(G)\)}{289}
\Idx{\(\underline{\mathrm{PL}}\)}{293}
\Idx{\(\underline{\mathrm{PLT}}\)}{293}
\Idx{带基根数据的 \(p\)-态射}{179}
\Idx{约化根数据的 \(p\)-态射}{100}
\Idx{根数据的权}{97}
\Idx{基本权}{97}
\Idx{处于一般位置的 Borel 子群}{158}
\Idx{两个抛物子群的相对位置}{296}
\Idx{两个抛物子群的标准位置}{308}
\Idx{两个抛物子群的横截位置}{297}
\Idx{\(\underline{\mathrm{PT}}\)}{293}
\Idx{拟可分裂的 \(S\)-约化群}{230}
\Idx{\(\underline{\mathrm{Q\acute ep}}\)}{232}
\Idx{拟钉扎}{230}
\Idx{约化群的中心商}{126}
\Idx{Chevalley 法则}{211}
\Idx{约化群概形的根}{13, 109}
\Idx{根数据的根}{64}
\Idx{无穷小根 \(\overline\alpha\)}{114}
\Idx{\(\rad^u(G)\)}{2}
\Idx{\(\rad(G)\)}{127}
\Idx{\(S\)-约化群的根基}{127}
\Idx{代数群的根基}{2}
\Idx{抛物子群的根基}{284}
\Idx{根数据的根基}{92}
\Idx{幺幂根基}{2}
\Idx{\textup{(RC)} 型子群的幺幂根基}{166, 281}
\Idx{抛物子群的幺幂根基}{281, 283}
\Idx{\(\operatorname{rad}(\mathscr R)\)}{92}
\Idx{\(\operatorname{rad}(G)\)}{2}
\Idx{\(\operatorname{rad}^u(P)\)}{281, 283}
\Idx{光滑仿射 \(k\)-群的约化秩}{4}
\Idx{根数据的约化秩}{64}
\Idx{光滑仿射 \(k\)-群的半单秩}{5}
\Idx{根数据的半单秩}{64}
\Idx{\(\operatorname{rgred}(G)\)}{4}
\Idx{\(\operatorname{rgred}(\mathscr R)\)}{64}
\Idx{\(\operatorname{rgss}(G)\)}{5}
\Idx{\(\operatorname{rgss}(\mathscr R)\)}{64}
\Idx{\(\mathscr R,\mathscr R^*\), 根数据}{64}
\Idx{\(\underline{\operatorname{Redext}}\)}{232}
\Idx{\(\underline{\operatorname{Rev}}\)}{232}
\Idx{\(\mathscr R(G)\)}{179}
\Idx{\(\mathcal R\), 根概形}{15}
\Idx{\(s_\alpha(t)\)}{110}
\Idx{约化群的 Dynkin 概形}{228}
\Idx{\(S\)-约化群的余根概形}{111}
\Idx{\(S\)-约化群的根概形}{15}
\Idx{约化群的抛物子群概形}{289}
\Idx{约化群的抛物子群的型概形}{291}
\Idx{约化群概形}{8}
\Idx{Hensel 局部概形}{220, 223, 261, 330}
\Idx{半局部概形}{162, 230, 239, 241, 287, 288, 311, 318, 322}
\Idx{\(\operatorname{sc}(\mathscr R)\)}{95}
\Idx{\(\operatorname{ss}(\mathscr R)\)}{95}
\Idx{\(s_G\)}{236}
\Idx{单连通 \(S\)-约化群}{127}
\Idx{单连通典范分解}{243}
\Idx{单连通根数据}{91}
\Idx{商群交换的子群}{174}
\Idx{Cartan 子群}{128}
\Idx{极小抛物子群}{316}
\Idx{Borel 子群}{131}
\Idx{分裂约化群的 Borel 子群}{139}
\Idx{Levi 子群}{282}
\Idx{\textup{(R)} 型子群}{131}
\Idx{具有约化纤维的 \textup{(R)} 型子群}{162}
\Idx{分裂约化群的 \textup{(R)} 型子群}{137}
\Idx{\textup{(RC)} 型子群}{165}
\Idx{抛物子群}{131, 279}
\Idx{临界约化子群}{163}
\Idx{\(\underline{\operatorname{Stand}}(G)\)}{309}
\Idx{根系}{64}
\Idx{约化群概形的根系}{14}
\Idx{正根系}{73}
\Idx{单根系}{71}
\Idx{Chevalley 系}{269}
\Idx{初等系}{28}
\Idx{初等系的生成元与关系}{58}
\Idx{Chevalley 系}{209}
\Idx{\(t_{\alpha\beta}=(w_\alpha w_\beta)^{n_{\alpha\beta}}\)}{180}
\Idx{\(T^{\mathrm{ad}}\)}{224}
\Idx{\(T_\alpha\)}{7}
\Idx{唯一性定理}{207}
\Idx{定理 90}{215}
\Idx{存在性定理}{267}
\Idx{Bruhat 定理}{149, 151}
\Idx{共轭定理}{311}
\Idx{基本定理}{202}
\Idx{\(S\)-约化群的余根基环面}{171}
\Idx{临界环面}{163}
\Idx{分裂环面}{10}
\Idx{\(S\)-群 \(G\) 的极大环面}{10, 23}
\Idx{已平凡化环面}{10}
\Idx{\(\underline{\operatorname{Tor}}(G)\)}{156}
\Idx{两个 \textup{(R)} 型子群的严格输运子}{134}
\Idx{抛物子群的严格输运子}{280}
\Idx{辫子关系}{212}
\Idx{约化群的型}{115}
\Idx{点 \(s\) 处的型}{115}
\Idx{\(\mathbf t\)}{289}
\Idx{\(\mathbf t(P)\)}{289}
\Idx{\(U_\alpha\)}{33, 110}
\Idx{\(V^\times\)}{16}
\Idx{\(w_\alpha\)}{7}
\Idx{\(w_\alpha(X)\)}{111}
\Idx{\(W(T)\)}{6}
\Idx{\(W(\mathscr F)^\alpha\)}{12}
\Idx{\(W(\mathscr R)\)}{65, 98}
\Idx{\(W^*(\mathscr R)=W(\mathscr R^*)\)}{65}
\Idx{\(W^*\)}{212}
\Idx{\(X_\alpha\)}{109}
\Idx{\(Z_\alpha\)}{7}
\Idx{\(Z(\mathscr R)\)}{98}
\Idx{\(Z^1(S'/S,G)\)}{262}

% END inputs/TERMINAL_INDEX/index/tex/components/00_tomeIII_index_entries.tex

\end{multicols}

% END inputs/TERMINAL_INDEX/index/tex/SGA3_TomeIII_Index_body.tex

